\documentclass[12pt,a4paper]{article}

\usepackage[utf8]{inputenc}
\usepackage[T2A]{fontenc}
\usepackage[serbian]{babel}

\usepackage{graphicx}
\usepackage{amsmath}
\usepackage{amsthm}
\usepackage{xcolor}
\usepackage{geometry}

\geometry{margin=2.5cm}

\newtheorem{definicija}{Дефиниција}

\title{Вештачка интелигенција у уметности}
\author{Тара Луциа Ревид}
\date{2026.}

\begin{document}

\renewcommand{\contentsname}{Садржај}

\maketitle

\tableofcontents

\section{Увод}

Вештачка интелигенција представља једну од најзначајнијих
технологија савременог доба. Њена примена постаје све присутнија
и у областима које су традиционално биле повезане са људском
креативношћу, као што су уметност, музика, фотографија и дизајн.

\section{Вештачка интелигенција и креативност}

Вештачка интелигенција може да анализира велики број података
и на основу њих генерише нови садржај.

\subsection{Примери примене}

Неки од примера употребе вештачке интелигенције у уметности су:

\begin{itemize}
    \item генерисање дигиталних слика,
    \item креирање музике,
    \item обрада фотографија,
    \item генерисање текста,
    \item обрада видео садржаја.
\end{itemize}

Редослед једног једноставног процеса може се приказати и као:

\begin{enumerate}
    \item унос података,
    \item анализа података,
    \item генерисање резултата.
\end{enumerate}

\begin{definicija}
Вештачка интелигенција је област информатике која се бави
развојем система способних да извршавају задатке који обично
захтевају људску интелигенцију.
\end{definicija}

\section{Математички пример}

Једна од основних формула која се користи у различитим
научним областима је:

\[
E = mc^2
\]

где је $E$ енергија, $m$ маса, а $c$ брзина светлости.

\section{Табеларни приказ}

\begin{center}
\begin{tabular}{|l|l|}
\hline
\textbf{Област} & \textbf{Пример примене AI} \\
\hline
Уметност & Генерисање слика \\
\hline
Музика & Генерисање мелодија \\
\hline
Видео & Монтажа и обрада \\
\hline
\end{tabular}
\end{center}

\section{Слика}

На следећој слици приказана је фотографија која је коришћена
и у HTML делу пројекта.

\begin{center}
    \includegraphics[width=0.35\textwidth]{slike/profil.jpg}
\end{center}

\section{Форматирање текста}

У документу се могу користити различити начини истицања текста.
Ово је \textbf{подебљан текст}, ово је \textit{курзив},
а ово је \textcolor{blue}{текст у боји}.

\section{Закључак}

Вештачка интелигенција има све значајнију улогу у савременом
креативном раду. Она може да буде користан алат за уметнике,
дизајнере и друге ауторе, али истовремено отвара питања
оригиналности, ауторства и односа између човека и технологије.

\end{document}